% Generated draft: requires mathematical and source review.
\documentclass[11pt]{article}
\usepackage[margin=25mm]{geometry}
\usepackage{fontspec}
\setmainfont{Latin Modern Roman}
\usepackage{amsmath,amssymb,mathtools,unicode-math}
\setmathfont{Latin Modern Math}
\usepackage{enumitem,needspace,xurl,hyperref}
\setlength{\parindent}{0pt}
\setlength{\parskip}{5pt}
\setlength{\emergencystretch}{3em}
\newcommand{\N}{\mathbb{N}}
\newcommand{\Z}{\mathbb{Z}}
\newcommand{\Q}{\mathbb{Q}}
\newcommand{\R}{\mathbb{R}}
\newcommand{\C}{\mathbb{C}}
\newcommand{\F}{\mathbb{F}}
\newcommand{\A}{\mathbb{A}}
\newcommand{\PP}{\mathbb{P}}
\newcommand{\E}{\mathbb{E}}
\newcommand{\eps}{\varepsilon}
\newcommand{\dd}{\,\mathrm{d}}
\begin{document}
\section*{UnsolvedMath ID: 701213}
\textbf{Source problem code:} catalogue entry
\subsection{Definitions and mathematical statement}
Let \(F\) be a field and let \(n\in\Z\) with \(n>2\). Write
\[
\mathrm{GL}_n(F)=\{S\in M_n(F):\det S\ne0\},\qquad
\mathrm{SL}_n(F)=\{A\in M_n(F):\det A=1\}.
\]
A matrix in \(\mathrm{SL}_n(F)\) is central if it belongs to
\[
Z(\mathrm{SL}_n(F))=\{\lambda I_n:\lambda\in F^{\times},\ \lambda^n=1\}.
\]
Use ordinary matrix similarity over \(F\): matrices \(A,B\in M_n(F)\) are similar if \(B=SAS^{-1}\) for some \(S\in\mathrm{GL}_n(F)\).

The question is whether, for every field \(F\), every integer \(n>2\), every \(\alpha\in\mathrm{SL}_n(F)\setminus Z(\mathrm{SL}_n(F))\), and every \(\beta\in\mathrm{SL}_n(F)\setminus Z(\mathrm{SL}_n(F))\), there exist \(S_1,\ldots,S_n\in\mathrm{GL}_n(F)\) such that
\[
\beta=(S_1\alpha S_1^{-1})(S_2\alpha S_2^{-1})\cdots(S_n\alpha S_n^{-1}).
\]
There are exactly \(n\) factors, with repetitions permitted. The source says ``similar'' without specifying the group of conjugating matrices; the formulation above uses its usual matrix-theoretic meaning. If conjugacy within \(\mathrm{SL}_n(F)\) was intended, that is a separate interpretation requiring clarification.
\subsection{Short English statement}
L. Vaserstein asks whether, over every field and for every \(n>2\), each noncentral determinant-one matrix is a product of exactly \(n\) matrices similar to any prescribed noncentral determinant-one matrix. Similarity here means conjugation by invertible matrices over the same field. The supplied registry labels the problem open; that status has not been independently verified.
\subsection{Sources}
\begin{itemize}
\item Local UnsolvedMath registry: supplied complete source record, permanent ID \(701213\), exact source problem code \texttt{catalogue entry}; upstream revision \texttt{d371ee917ac9647fbb06cd6fc529deec6f3dda82}. The record was created and updated on \texttt{2026-10-16T14:00:00Z} and supplies the status ``open.''
\item The supplied attribution is to L. Vaserstein, proposed in 2006, in \emph{The Kourovka Notebook (Unsolved Problems in Group Theory)}, Problem 16.21, 16th issue (2006). The background also identifies the 21st issue edition (2026), \texttt{arXiv:1401.0300v48}. These bibliographic details are supplied source information, not independently verified.
\item The supplied collection description identifies E. I. Khukhro and V. D. Mazurov as editors of the Kourovka Notebook collection covering Issues 1--21 (1965--2026).
\end{itemize}
\end{document}
